%BeginMC
\question
Early successional species add nutrients and deepen the soil, allowing intermediate successional species to colonize an area. This process is known as

\begin{choices}
\correctchoice facilitation. %EndChoice
\choice tolerance. %EndChoice
\choice inhibition. %EndChoice
\choice secondary succession. %EndChoice
\choice primary succession. %EndChoice
\choice disturbance. %EndChoice
\end{choices}
%EndMC

%BeginMC
\question
Of the following choices, which is the biggest concern about climate change?

\begin{choices}
\correctchoice The rate of temperature increase is faster than ever recorded. %EndChoice
\choice The amount of solar radiation will change faster than plants can adapt. %EndChoice
\choice Forest and grassland ecosystems will die due to less \ce{CO2} available for photosynthesis. %EndChoice
\choice More regions around the globe will experience stronger rain shadows. %EndChoice
\choice More regions around the globe will experience desert-like conditions. %EndChoice
\choice The average global temperature will be warmer than ever recorded.  %EndChoice
\end{choices}
%EndMC

%BeginMC
\question
What type of survivorship curve describes a species that has a high rate of juvenile mortality but low adult mortality?

\begin{choices}
\choice Type I curve. %dontshuffle %EndChoice
\choice Type II curve. %dontshuffle %EndChoice
\correctchoice Type III curve. %dontshuffle %EndChoice
\choice Type IV curve %dontshuffle %EndChoice
\end{choices}
%EndMC


%BeginMC
\question
In a population growing according to the exponential growth model, population size is

\begin{choices}
\correctchoice not limited. %EndChoice
\choice limited by density-dependent factors. %EndChoice
\choice limited by density-independent factors. %EndChoice
\choice limited by both density-dependent and density-independent factors. %EndChoice
\end{choices}
%EndMC


%BeginMC
\question
An organism that only reproduces once is known as

\begin{choices}
\correctchoice semelparous. %EndChoice
\choice iteroparous. %EndChoice
\choice ectotherm. %EndChoice
\choice perennial. %EndChoice
\end{choices}
%EndMC

%BeginMC
\question
Which dispersion should be favored among plants that compete intraspecifically by casting shade?

\begin{choices}
\correctchoice Evenly spaced. %EndChoice
\choice Clumped. %EndChoice
\choice Random. %EndChoice
\choice Clustered. %EndChoice
\end{choices}
%EndMC


%BeginMC
\question
Probably the most important factor(s) affecting the global distribution of ecosystems is (are)

\begin{choices}
\correctchoice climate.%EndChoice
\choice wind and ocean water current patterns.%EndChoice
\choice wind direction and rain shadows.%EndChoice
\choice longitude and latitude.%EndChoice
\choice day length and rainfall.%EndChoice
\end{choices}
%EndMC

%BeginMC
\question
According to this climograph, the temperate grassland 

\begin{multicols}{2}
	\begin{choices}
		\choice never occurs with deserts. %EndChoice
		\choice has an average annual rainfall of about 20--80 cm. %EndChoice
		\choice has an average annual rainfall of about 5--30 cm. %EndChoice
		\correctchoice has an average annual temperature range of about 4--30°C. %EndChoice
		\choice has the greatest variation in precipitation of any ecosystem shown in the figure. %EndChoice
	\end{choices}
\columnbreak
	{\centering \includegraphics[width=0.9\columnwidth]{exam4_climograph1}\par}
\end{multicols}
%EndMC

%BeginMC
\question
Organisms that break down organic material and make it available to other organisms are called

\begin{choices}
\correctchoice decomposers. %EndChoice
\choice scavengers. %EndChoice
\choice consumers.  %EndChoice
\choice producers. %EndChoice
\choice degraders. %EndChoice
\end{choices}
%EndMC

%BeginMC
\question
Energy from the sun is converted into chemical energy for organisms by

\begin{choices}
\correctchoice primary producers.%EndChoice
\choice primary consumers.%EndChoice
\choice carnivores.%EndChoice
\choice decomposers.%EndChoice
\choice herbivores.%EndChoice
\end{choices}
%EndMC

%BeginMC
\question
The energy that is \textit{not} transferred to other trophic levels or decomposers is lost

\begin{choices}
\correctchoice as metabolic heat.%EndChoice
\choice in fecal matter.%EndChoice
\choice to consumers at any trophic level.%EndChoice
\choice to primary consumers.%EndChoice
\choice to secondary consumers.%EndChoice
\end{choices}
%EndMC


%BeginMC
\question
A field containing 2000 kg of plant material can support about \rule{1in}{0.4pt} of tertiary consumers.

\begin{choices}
\correctchoice 2 kg %dontshuffle %EndChoice
\choice 10\% %dontshuffle %EndChoice
\choice 10 kg %dontshuffle %EndChoice
\choice 20 kg %dontshuffle %EndChoice
\choice 90\% %dontshuffle %EndChoice
\choice 200 kg %dontshuffle %EndChoice
\end{choices}
%EndMC


%BeginMC
\question
Based on the hypothetical food web below, one letter represents an organism that could be both a secondary and tertiary consumer? The arrows point in the direction of energy transfer from lower to higher trophic levels. More than one answer is possible. Circle only one answer.

\begin{multicols}{3} 
	\begin{choices}
		\correctchoice A %dontshuffle %EndChoice
		\choice B %dontshuffle %EndChoice
		\choice C %dontshuffle %EndChoice
		\choice D %dontshuffle %EndChoice
		\correctchoice E %dontshuffle %EndChoice
	\end{choices}
\columnbreak
	\includegraphics{exam4_foodweb}
\columnbreak
\end{multicols}
\vspace{-1\baselineskip}
%EndMC

%BeginMC
\question
The ``role'' or ``job'' of a species in a community, represented by all of the resources needed by the species, is called 

\begin{choices}
\correctchoice the niche.%EndChoice
\choice the resource.%EndChoice
\choice a predator.%EndChoice
\choice a predator or prey.%EndChoice
\choice a producer or consumer.%EndChoice
\end{choices}
%EndMC

%BeginMC
\question
White-breasted nuthatches and brown creepers both eat insects that hide in the furrows of bark in hardwood trees. The brown creeper searches for insects by hunting from the bottom of the tree trunk toward the top, whereas the white-breasted nuthatch searches from the top of the trunk down. Their feeding areas rarely overlap on the tree trunk. These hunting behaviors best illustrate which of the following ecological concepts?

\begin{choices}
\choice competitive exclusion%EndChoice
\correctchoice resource partitioning%EndChoice
\choice intraspecific competition%EndChoice
\choice keystone species%EndChoice
\choice interspecific competition%EndChoice
\end{choices}
%EndMC



%BeginMC
\question
Competitive exclusion means that

\begin{choices}
\correctchoice one species will always out complete another species that occupies the same niche.  %EndChoice
\choice two species will partition resources to eliminate competition. %EndChoice
\choice one individual will always out complete another individual in a population. %EndChoice
\choice one population will always exclude another another population of the same species in a community. %EndChoice
\choice a species will have a smaller realized niche than a fundamental niche. %EndChoice
\end{choices}
%EndMC

%BeginMC
\question
One form of elephantiasis occurs when \textit{Brugia malayyi}, a thread-like worm, invades human lymph nodes. The worm blocks the flow of lymph throughout the body, causing massive swelling of the lower body, especially legs and genitals. As described here, this is an example of

\begin{choices}
\choice commensalism.%EndChoice
\correctchoice parasitism.%EndChoice
\choice mutualism.%EndChoice
\choice predation.%EndChoice
\choice symbiosis.%EndChoice
\end{choices}
%EndMC


%BeginMC
\question
Symbiosis is a long-term ecological interaction

\begin{choices}
\correctchoice between two or more species. %EndChoice
\choice where one species benefits and another species is harmed %EndChoice
\choice where both species benefit. %EndChoice
\choice  where one species benefits and the other species neither benefits nor is harmed. %EndChoice
\choice where two or more species partition resources. %EndChoice
\end{choices}
%EndMC

%BeginMC
\question
Colors and patterns that helps an animal blend in with its environment is called

\begin{choices}
\correctchoice cryptic coloration. %EndChoice
\choice environmental mimicry. %EndChoice
\choice warning coloration %EndChoice
\choice aposomatic coloration. %EndChoice
\choice habitat mimicry. %EndChoice
\end{choices}
%EndMC

%BeginMC
\question
Which block in the figure below shows the greatest species diversity? (Different shades, including white, represent different species. The small squares indicate relative abundance of individuals present.)

\begin{multicols}{3}
\begin{choices}
\correctchoice \includegraphics{diversity_mcE} %EndChoice
\choice \includegraphics{diversity_mcA} %EndChoice
\choice \includegraphics{diversity_mcB} %EndChoice
\choice \includegraphics{diversity_mcC} %EndChoice
\choice \includegraphics{diversity_mcD} %EndChoice
\end{choices}
\end{multicols}
%EndMC


%BeginMC
\question
Fire suppression by humans

\begin{choices}
\choice will always result in an increase in species diversity in a given biome. %EndChoice
\correctchoice can change the species composition within biological communities. %EndChoice
\choice will result ultimately in sustainable production of increased amounts of forest products for human use. %EndChoice
\choice is necessary for the protection of threatened and endangered forest species. %EndChoice
\choice is a management goal of conservation biologists to maintain the healthy condition of forest communities.%EndChoice
\end{choices}
%EndMC

%BeginMC
\question
The difference between communities and ecosystems is that

\begin{choices}
\choice communities include individuals of the same species; ecosystems include individuals of different species.%EndChoice
\choice communities include individuals of the different species; ecosystems are individuals of the same species.%EndChoice
\choice communities include interactions between individuals and their physical environment; ecosystems do \emph{not} include such interactions.%EndChoice
\correctchoice communities do \emph{not} include interactions between individuals and the physical environment; ecosystems \emph{do} include such interactions.%EndChoice
\choice communities are groups of individuals of the different species; ecosystems are groups of communities.%EndChoice
\end{choices}
%EndMC

%BeginMC
\question
Plants produce water vapor during photosynthesis, which goes from the plant to the atmosphere. The water vapor in the atmosphere eventually falls back to the earth as rain. This process would be part of

\begin{choices}
\correctchoice the water cycle.%EndChoice
\choice energy flow.%EndChoice
\choice plant metabolism.%EndChoice
\choice the entire ecosystem.%EndChoice
\choice the energy cycle.%EndChoice
\choice the nutrient cycle.%EndChoice
\end{choices}
%EndMC

%BeginMC
\question
Which statement most accurately describes how nutrients and energy are used in ecosystems?

\begin{choices}
\choice Nutrients flow through ecosystems; energy cycles within ecosystems. %EndChoice
\correctchoice Energy flows through ecosystems; nutrients cycle within ecosystems. %EndChoice
\choice Energy can be converted into nutrients; nutrients cannot be converted into energy. %EndChoice
\choice Nutrients can be converted into energy; energy cannot be converted into nutrients. %EndChoice
\choice Nutrients are used in ecosystems; energy is not. %EndChoice
\end{choices}
%EndMC


%BeginMC
\question
Which of the following statements about ecosystems is false?

\begin{choices}
\correctchoice On average, 10\% of the energy available at one trophic is lost as it moves to the next higher trophic level. %EndChoice
\choice Energy flows through an ecosystem. %EndChoice
\choice Nutrients, like nitrogen and phosphorus, cycle within an ecosystem. %EndChoice
\choice Water cycles within an ecosystem. %EndChoice
\choice Carnivores occupy higher trophic levels than herbivores. %EndChoice
\end{choices}
%EndMC

